\subsection{Early cut-off}

We have a deterministic trajectory for the counter and a concentration estimate while a positive fraction of
the pieces remain. We now
choose the cut-off so that the queens placed afterward can destroy only a fraction of the repairs already
available. For this purpose,
fix \begin{equation}\label{eq:constantes}
 \eta=\kappa/2,\qquad a_0=\eta^4/16,\qquad \delta_0=a_0/16;
\end{equation} so that $a_0>8\delta_0$ and $\delta_0<1/4$. For a cut-off $\delta\in(0,1/4)$, let
$m_\delta=\lceil\delta n\rceil$, let
$\tau_\delta$ be the time when exactly $m_\delta$ real pieces remain to be placed, and let $\zeta=\delta/4$.
We use the cut-off
$\delta_0$ in the lower bound (Section~\ref{sec:positiva}). In Sections~\ref{sec:inversa}
and~\ref{sec:limite} the cut-off $\delta$ is
a free parameter and is sent to zero.

\begin{lema}[Absorbers at the cut-off]\label{lem:corte} Let $h\in\{1,2\}$ and $\delta=\delta_0$. Except on
an event of probability
$e^{-cn}$ under $\QQ_h$, the process reaches the cut-off with $\tau_{\delta_0}\le T_{\delta_0/4}$ and,
simultaneously for all roots,
\begin{equation}\label{eq:corte}
 A^{rc}_{\tau_{\delta_0}}\ge a_0n .
\end{equation} \end{lema}

\begin{proof} Concentration of the colors in~\eqref{eq:regularidad}, with tolerance smaller than
$\delta_0/8$, ensures
$\tau_{\delta_0}\le T_{\delta_0/4}$ except on an event of probability $e^{-cn}$. Deterministically, at the
cut-off
$y=\tau_{\delta_0}/N\ge(q-m_{\delta_0})/N\ge1/4$ for $n$ large. Moreover,
$\sum_{e\in\mathcal I_{rc}}\pi_e\ge q-O(1)\ge n/2$, since
$\sum_e\pi_e=q$ and the excluded squares number $O(n)$ with $\pi_e\le C/n$. By Lemma~\ref{lem:balance},
$1-yp_{d_i}\ge1-p_{d_i}\ge\eta$, so $a_n^{rc}(y)\ge\eta^4n/8$ for all $y\ge1/4$.
Lemma~\ref{lem:conc_directa} with tolerance
$\gamma=\eta^4/16$, applied to all deterministic horizons $T\le T_{\delta_0/4}$ with $T/N\ge1/4$ (a union
bound over $O(n)$ horizons
allows for the random cut-off), gives $A^{rc}_{\tau_{\delta_0}}\ge\eta^4n/16=a_0n$ for all roots. The claim
only requires the process
to reach the cut-off. Lemma~\ref{lem:regularidad} already accounts for any earlier failure. \end{proof}

\emph{Order of choices.} Fix $\eta,a_0,\delta_0$ and $\zeta=\delta_0/4$, then choose the tolerances in
Lemmas~\ref{lem:regularidad}--\ref{lem:conc_directa}, which depend on these parameters, and finally take $n$
large. Uniformity as
$\delta\to0$ is not required. The constants come from analytic bounds and are very conservative.

\subsection{Deterministic loss after the cut-off}

\begin{lema}\label{lem:perdida} Let $\delta\in(0,1/4)$ and let $(r,c)$ be a root whose row and column remain
physically empty. Each
real piece creates at most one occupied candidate for $(r,c)$ and cancels the contributions of at most four.
Virtual pieces leave
$A^{rc}$ unchanged. Consequently, on the success event, for every crossing that is empty at the end, \[
 A^{rc}_{\mathrm{final}}\ge A^{rc}_{\tau_\delta}-4m_\delta,\qquad
 |A^{rc}_{\mathrm{final}}-A^{rc}_{\tau_\delta}|\le4m_\delta .
\] \end{lema}

\begin{proof} A new queen at square $g$ can create only the term $Z^\oc_{g,\varnothing}$. It destroys the
term for an occupied
candidate $e$ if $g$ lies on one of the destination diagonals of $e$. The square $g$ lies on one diagonal
and one antidiagonal, and each is
a destination diagonal for at most two occupied candidates, since each row and column contains at most one
queen. Thus each real
piece changes $A^{rc}$ by an integer in $[-4,1]$. Virtual pieces merely
deactivate vertices, and $A^{rc}$ depends only on $O$ and $V$. Exactly $m_\delta$ real pieces are placed
after the cut-off. \end{proof}
